\section*{Capítulo \ref{ch:b3:deformation-mechanisms} --- Mecanismos de deformação e zonas de cisalhamento}

\begin{solution}{exo:b3:deformation-mechanisms:1}
Forsterita: $1573/2163 = 0.73$; gelo: $263/273 = 0.96$. O gelo está muito mais próximo
de seu ponto de fusão e flui com mais facilidade: as geleiras escoam na superfície
da Terra.
\end{solution}

\begin{solution}{exo:b3:deformation-mechanisms:2}
$\gamma = \tan\ang{70} = 2.75$ e $\tan\ang{80} = 5.67$. Próximo de \ang{90}, a
tangente varia muito rapidamente com o ângulo: um grau de erro em $\psi =
\ang{84}$ altera $\gamma$ em um quinto.
\end{solution}

\begin{solution}{exo:b3:deformation-mechanisms:3}
$\sigma = (10^{3.56}/10)^{1/1.26} = 363.1^{0.794} = \qty{108}{MPa}$.
\end{solution}

\begin{solution}{exo:b3:deformation-mechanisms:4}
$E/RT = 223000/(8.314\times773.15) = 34.69$, logo $\dot\varepsilon =
\num{1.1e-4}\times20^4\times e^{-34.69} = \num{1.1e-4}\times\num{1.6e5}\times
\num{8.6e-16} = \qty{1.5e-14}{s^{-1}}$.
\end{solution}

\begin{solution}{exo:b3:deformation-mechanisms:5}
$\eta = \sigma/(2\dot\varepsilon)$: $\num{65.4e6}/\num{2e-14} =
\qty{3.3e21}{Pa.s}$ a \qty{400}{\celsius} e $\num{18.0e6}/\num{2e-14} =
\qty{9.0e20}{Pa.s}$ a \qty{500}{\celsius}: a crosta média, nessas
temperaturas, tem viscosidade aproximadamente igual à do manto abaixo dos mantos de gelo.
\end{solution}

\begin{solution}{exo:b3:deformation-mechanisms:6}
$n = \log(\num{1.6e-4}/\num{e-5})/\log2 = \log16/\log2 = 4$: \omterm{def:b3:deformation-mechanisms:creep}{fluência por deslocações},
cujo \omterm{def:b3:deformation-mechanisms:flow-law}{expoente de tensão} é bem superior a 1.
\end{solution}

\begin{solution}{exo:b3:deformation-mechanisms:7}
$\ln10 = (E/R)(1/1000 - 1/1100)$, logo $E = 8.314\times2.3026/\num{9.091e-5} =
\qty{211}{kJ.mol^{-1}}$.
\end{solution}

\begin{solution}{exo:b3:deformation-mechanisms:8}
Nabarro--Herring ($d^{-2}$): 100 vezes mais rápida; Coble ($d^{-3}$): 1000 vezes.
Ao longo dos contornos, os átomos que se difundem usam um caminho cuja seção transversal, $\delta
d$, diminui com o grão, além da menor distância e do menor
volume a deslocar.
\end{solution}

\begin{solution}{exo:b3:deformation-mechanisms:9}
$\sigma_b^{2.5} = \num{3.59e-6}\,d^{-3}/\num{6.54e-17}$: para $d = 1000$,
$\sigma_b = 54.9^{0.4} = \qty{5.0}{MPa}$; para $d = 100$, $\sigma_b =
\num{54900}^{0.4} = \qty{79}{MPa}$. A inclinação é $-3/2.5 = -1.2$: um grão
dez vezes menor eleva o contorno em $10^{1.2} = 16$.
\end{solution}

\begin{solution}{exo:b3:deformation-mechanisms:10}
Nos dois caminhos, a diferença de concentração motriz é a mesma e a
distância é $d$; o volume a deslocar por unidade de deformação é $d^3$. A rede
oferece uma seção transversal $\propto d^2$, logo a taxa varia como $d^2/(d\,d^3) =
d^{-2}$; os contornos oferecem $\propto\delta d$, logo $\delta d/(d\,d^3) = \delta
d^{-3}$. A razão entre as taxas de contorno e de rede é $(B_B\delta D_B)/(B_LD_Ld)$,
maior que 1 quando $d < (B_B/B_L)\,\delta D_B/D_L$.
\end{solution}

\begin{solution}{exo:b3:deformation-mechanisms:11}
$\eta \propto \dot\varepsilon^{(1-n)/n} = \dot\varepsilon^{-2.5/3.5}$: dez vezes
mais rápido divide a viscosidade por $10^{0.714} = 5.2$. Uma faixa que, por acaso,
se deforma mais rápido torna-se mais fraca, e então absorve mais da deformação:
a deformação concentra-se nela.
\end{solution}

\begin{solution}{exo:b3:deformation-mechanisms:12}
Antes: $\dot\varepsilon_{\mathrm{dis}} = \num{6.54e-17}\times10^{3.5} =
\qty{2.07e-13}{s^{-1}}$ e $\dot\varepsilon_{\mathrm{dif}} = \num{3.59e-6}
\times10/3000^3 = \qty{1.3e-15}{s^{-1}}$: \qty{2.08e-13}{s^{-1}}, \omterm{def:b3:deformation-mechanisms:creep}{fluência por
deslocações}. Depois: $\dot\varepsilon_{\mathrm{dif}} = \num{3.59e-6}\times10/30^3 =
\qty{1.33e-9}{s^{-1}}$, \omterm{def:b3:deformation-mechanisms:creep}{fluência por difusão}, \num{6400} vezes mais rápida sob a
mesma tensão: a recristalização amoleceu enormemente a zona.
\end{solution}

\begin{solution}{pb:b3:deformation-mechanisms:1}
\textbf{1.} O dique passa de \ang{90} para \ang{6} em relação à zona: $\psi =
\ang{84}$ e $\gamma = \tan\ang{84} = 9.51$.

\textbf{2.} $\gamma w = 9.51\times400 = \qty{3804}{m}$, cerca de \qty{3.8}{km}.

\textbf{3.} $S_1 = \sqrt{1 + 9.51^2/4} + 9.51/2 = 4.859 + 4.755 = 9.614$,
$S_2 = 1/S_1 = 0.104$, e $S_1/S_2 = 92.4$.

\textbf{4.} $\tan2\theta' = 2/9.51 = 0.210$, $\theta' = \ang{5.9}$: a
foliação, próxima do eixo longo, gira cada vez mais para perto da zona à medida
que $\gamma$ cresce.

\textbf{5.} Seguindo o dique para o norte, ele é arrastado para leste dentro da zona e
deixa-a mais a leste: a parede norte deslocou-se para leste em relação à parede sul.
Visto da parede sul, o lado distante deslocou-se para a direita: a zona é
dextral.

\textbf{6.} No plano que contém a lineação de estiramento, horizontal aqui
por se tratar de uma zona transcorrente, e perpendicular à foliação vertical: uma
seção horizontal. Sua foliação S inclina-se no sentido anti-horário em relação à zona,
em direção a lés-nordeste, como o eixo longo da elipse de deformação.

\textbf{7.} $\sigma = (10^{3.56}/20)^{1/1.26} = 181.5^{0.794} = \qty{62}{MPa}$.

\textbf{8.} $242.1^{0.794} = \qty{78}{MPa}$ para \qty{15}{\micro m} e
$145.2^{0.794} = \qty{52}{MPa}$ para \qty{25}{\micro m}.

\textbf{9.} Sim: os experimentos abrangeram de 34 a \qty{268}{MPa}.

\textbf{10.} A zona deformou-se a \qty{450}{\celsius}, muito abaixo dos 800 a
\qty{1100}{\celsius} dos experimentos; uma relação que não depende da
temperatura pode ser estendida até ela, embora ainda seja uma
extrapolação.

\textbf{11.} $D = 3631/100^{1.26} = 3631/331 = \qty{11}{\micro m}$.

\textbf{12.} $E/RT = 223000/(8.314\times723.15) = 37.09$ e $e^{-37.09} =
\num{7.79e-17}$.

\textbf{13.} $\dot\varepsilon = \num{1.1e-4}\times62^4\times\num{7.79e-17} =
\num{1.1e-4}\times\num{1.478e7}\times\num{7.79e-17} = \qty{1.27e-13}{s^{-1}}$.

\textbf{14.} $\dot\gamma = 2\dot\varepsilon = \qty{2.53e-13}{s^{-1}}$.

\textbf{15.} $\eta = \num{62e6}/(2\times\num{1.27e-13}) = \qty{2.4e20}{Pa.s}$,
próxima de \qty{3e20}{Pa.s} do manto superior sob o Lago Ardo: a zona
era tão fraca quanto o manto.

\textbf{16.} $e^{-E/RT}$ muda pelo fator 0.064 a \qty{400}{\celsius}
(dezesseis vezes mais lento) e 11.0 a \qty{500}{\celsius}.

\textbf{17.} $(78/62)^4 = 2.5$ vezes mais rápido.

\textbf{18.} $t = 9.51/\num{2.53e-13} = \qty{3.75e13}{s}$.

\textbf{19.} $\num{3.75e13}/\num{3.156e10} = \qty{1190}{kyr}$.

\textbf{20.} $\qty{3804}{m}/\qty{1.19e6}{yr} = \qty{3.2}{mm/yr}$.

\textbf{21.} Dez ou mais vezes mais lenta do que os limites principais de placas: a zona
era um ramo secundário, ou um de vários que compartilhavam o movimento de um limite.

\textbf{22.} $\dot\gamma = 0.02/(\num{3.156e7}\times400) = \qty{1.58e-12}{s^{-1}}$,
6.2 vezes mais rápido; a tensão sobe por $6.2^{1/4} = 1.58$, para \qty{98}{MPa}, e
os grãos teriam $3631/98^{1.26} = \qty{11}{\micro m}$.

\textbf{23.} O tempo varia como $\sigma^{-4}$: a faixa de tamanho de grão, sozinha, fornece
$(62/78)^4 = 0.40$ a $(62/52)^4 = 2.0$ vezes a estimativa, e a
incerteza do expoente $1.26 \pm 0.13$ e das constantes da \omterm{def:b3:deformation-mechanisms:flow-law}{lei de fluxo}
amplia-a ainda mais: o resultado é uma ordem de grandeza, não uma data.

\textbf{24.} A \omterm{def:b3:deformation-mechanisms:shear-zone}{zona de cisalhamento} Varden levou cerca de
\textbf{\qty[text-series-to-math=true]{1190}{kyr}} para acumular sua deformação.
\end{solution}
